\chapter{Secularism, Pluralism, Fundamentalism \& Politics and Society}


\section{Secularism \& Pluralism}

\subsection{Indian vs Western (American) Secularism}

\begin{KeyConcept}
\textbf{The two core principles (GS framing):}
\begin{itemize}
\item \textbf{Equal respect for all religions}, and \textbf{non-interference of the state} in religious practices --- these are the two core principles. But note: these two are the \textbf{Indian model} of secularism, not secularism as such (one cannot define secularism without qualifying it as Indian/American/French).
\end{itemize}
\end{KeyConcept}

\begin{ComparisonBox}
\textbf{American vs Indian secularism (five differences):}
\begin{enumerate}
\item \textbf{Mutual exclusion (America) vs principled distance (India)} --- America runs on \textbf{mutual exclusion} (strict non-interference of religion and state in each other's affairs); India runs on \textbf{equidistance and principled distance} (distance based on principles; if the principles are violated, the state comes closer). ``Principled distance'' was coined by sociologist \textbf{Rajeev Bhargava}.
\item \textbf{Inter- and intra-religious domination} --- America does not attend to intra-religious domination (e.g., upper caste over lower caste within Hinduism; men over women via sati) or inter-religious domination (Hindu oppressing Muslim on religious grounds). In India, \textbf{secularism is a response to both} inter- and intra-religious domination.
\item \textbf{Manager vs trustee} --- the American state cannot manage any religious body; the Indian state is a \textbf{trustee/manager} of religious organisations (religious endowment departments, waqf boards, appointment of trustees).
\item \textbf{Reformism} --- the American state cannot reform religion; the Indian state is a \textbf{reformist, welfare state} (justice is social, not just economic --- it can equally ``disrespect'' a practice; e.g., triple talaq).
\item \textbf{Individualism vs collectivism} --- America gives little scope for community/minority rights (e.g., no Jamia-like institution); India's secularism covers the religious freedom of \textbf{individuals and minority communities} (Article 30(1) minority institutions).
\end{enumerate}
\end{ComparisonBox}

\begin{ExampleBox}
\begin{itemize}
\item The Jagannath-temple example: in India, I can go to court if denied entry (fundamental rights); in America the court would not intervene (religion is private). Sati: in India the court would strike it (violative of women's fundamental rights); in America the court treats it only as murder, not a religious matter.
\item \textbf{Majid denied accommodation for being Muslim} --- in India he can go to court (denial on religious grounds); in America he cannot.
\end{itemize}
\end{ExampleBox}

\begin{CriticalPoint}
\textbf{America vs India historically:}
\begin{itemize}
\item America has been historically \textbf{homogenous} (hence individual rights); India has been \textbf{heterogeneous since forever} (land of Buddha, Guru Nanak, Mahavira --- hence community-based rights). Rajeev Bhargava suggests America today \textbf{needs} principled distance more than India (being a major victim of terrorism).
\end{itemize}
\end{CriticalPoint}

\begin{DataFact}
\textbf{India's inherent secularism:}
\begin{itemize}
\item \textbf{Kabir} (``God is neither in Kailash, nor in idols, nor in the mosque --- all are one''); \textbf{Gandhi} (``Ishwar Allah tero naam''); \textbf{Nehru} (\textit{Glimpses of World History} --- Ashoka and Akbar spoke in the same tenor, the voice of India); \textbf{Akbar's Ibadat Khana} (a court for inter-faith dialogue, Sulh-i-kul / universal peace); \textbf{Ashoka's policy of Dhamma} and the \textbf{Seventh Major Pillar Edict} (``I wish all the pasandas cohabit my kingdom'') --- religious tolerance.
\end{itemize}
\end{DataFact}

\begin{QuickRevision}
\begin{itemize}
\item Two principles (equal respect + non-interference) = Indian model, not secularism as such.
\item America: mutual exclusion, private religion, individualist, no reform/management, blind to inter/intra-religious domination. India: principled distance, reformist state, trustee of religions, collective rights, response to inter/intra domination.
\item Inherent Indian secularism: Kabir, Gandhi, Ashoka (Dhamma, 7th edict), Akbar (Ibadat Khana), Nehru.
\end{itemize}
\end{QuickRevision}

\sectiondivider

\subsection{T. N. Madan}

\begin{CriticalPoint}
\textbf{The two camps of scholars:}
\begin{itemize}
\item \textbf{Against imposing Western secularism on India (anti-Western secularism):} \textbf{T. N. Madan, Ashis Nandy, Partha Chatterjee} (they discuss the \textit{failure} of secularism in India). \textbf{Pro-secularism:} \textbf{Rajeev Bhargava} and \textbf{Amartya Sen} (and the Western \textbf{D. Roover}) --- they favour a unique, neither-Western-nor-Indian secularism.
\end{itemize}
\end{CriticalPoint}

\begin{DataFact}
\textbf{T. N. Madan's argument:}
\begin{itemize}
\item Western secularism (strict separation) cannot simply be copy-pasted onto Indian soil, because \textbf{religion is in the blood of Indians} --- it shapes food behaviour, clothing, rituals, birth, marriage and death customs. Religion is a \textbf{way of life}; how can a way of life be separated from the state?
\item In the West religion can be kept wholly private (even religious public holidays exist so the state doesn't display religion; religion appears only in private/convent schools, not government schools) --- but this is impossible in India.
\item There is a \textbf{correlation between secularization and communal/sectarian conflict in India}: trying to become secular like America/France produces communal conflict.
\item Conclusion: for Western secularism to succeed in India, \textbf{religion must be given equal importance alongside it} --- but Western secularism is against religion-in-public, so the two are opposite.
\end{itemize}
\end{DataFact}

\begin{QuickRevision}
\begin{itemize}
\item T. N. Madan: religion is in Indians' blood (a way of life); strict separation impossible; secularization correlates with communal conflict; religion and Western secularism are opposite.
\end{itemize}
\end{QuickRevision}

\sectiondivider

\subsection{Ashis Nandy}

\begin{CriticalPoint}
\textbf{Religion pushed out comes back through the back door:}
\begin{itemize}
\item If you push religion out of personal life, it \textbf{returns through the back door in a violent form}.
\end{itemize}
\end{CriticalPoint}

\begin{KeyConcept}
\textbf{Two types of religion:}
\begin{itemize}
\item \textbf{Religion as faith} (belief) and \textbf{religion as ideology} (communalist political identity). With \textbf{modernization}, religion-as-faith \textit{may decline} (individualism, a man-centric world), but \textbf{religion-as-ideology accelerates} --- because political parties are a modern phenomenon, and we now have Hindu/Muslim religious political parties contesting each other. This is a \textbf{modern, sophisticated form of communalism} (communalization of politics).
\end{itemize}
\end{KeyConcept}

\begin{KeyConcept}
\textbf{Two forms of secularism:}
\begin{itemize}
\item \textbf{Western secularism} keeps religion out of politics (mutual exclusion); but \textbf{strict separation makes religion re-enter public life through the back door $\rightarrow$ communalization of politics}; Western secularism \textbf{alienates believers and breeds old and new violence} (old = fanaticism; new = fault-line conflict between secular and religious nation-states, e.g., America vs Middle East --- Samuel Huntington's \textbf{clash of civilizations}; the collapse of the USSR saw nation-states form on religious lines).
\item \textbf{Non-Western secularism} (he avoids ``Indian'' because India never had that separation): \textbf{continuous dialogue between faith-holders, bracketed by tolerance} --- promoted by Ashoka, Akbar, Gandhi. Indian secularism came from Buddhism, Islam and Hinduism --- not from America's secular politics.
\end{itemize}
\end{KeyConcept}

\begin{QuickRevision}
\begin{itemize}
\item Nandy: religion pushed out $\rightarrow$ returns violently; two religions (faith declines, ideology accelerates with modernization); Western secularism keeps religion out (fails, alienates, breeds violence) vs non-Western = continuous dialogue/tolerance (Ashoka, Akbar, Gandhi).
\end{itemize}
\end{QuickRevision}

\sectiondivider

\subsection{Partha Chatterjee}

\begin{CriticalPoint}
\textbf{The Indian Marxist's inherent contradictions:}
\begin{itemize}
\item Since the birth of Indian civilization, the Indian Marxist has been a victim of inherent contradictions. Secularism implies state--religion distance, yet the state also intervenes in religious affairs --- a contradiction (first ``I won't interfere'', then it intervenes, then again not --- hide-and-seek). The state--religion relation is a \textbf{moving equilibrium, not static}.
\end{itemize}
\end{CriticalPoint}

\begin{DataFact}
\textbf{Chatterjee's observations:}
\begin{itemize}
\item Secularism implied separation, but the state's \textbf{reformist intervention has been largely in the religious sphere of Hindus} (sati, widow-burning, animal sacrifice, untouchability, Hindu Marriage Act --- polygamy made illegal for Hindus but not Muslims, right to divorce introduced, child marriage abolished, inter-caste marriage recognised). So it deserves the \textbf{charge of favouritism and minoritism} (state changed Hindu personal laws significantly; Muslim divorce rights came very late, after the abolition of triple talaq).
\item He writes in the context of \textbf{Congress rule}; he is only recording, not value-loading. He concludes: India would not be secular unless it is \textbf{equally interventionist in all religions} --- the Indian model (equal respect $\rightarrow$ equal disrespect) is \textit{more} disrespectful towards Hindus, less towards Muslims.
\item (Note: SC status/reservation is given only to Hindu, Buddhist and Sikh SCs --- not Muslim/Christian SCs --- a selective intervention.)
\end{itemize}
\end{DataFact}

\begin{QuickRevision}
\begin{itemize}
\item Chatterjee: state intervenes selectively, mostly in Hindu practices $\rightarrow$ favouritism/minoritism; state--religion is a moving equilibrium; India is not secular unless equally interventionist in all religions.
\end{itemize}
\end{QuickRevision}

\sectiondivider

\subsection{Rajeev Bhargava: Two Types of Secularism}

\begin{KeyConcept}
\textbf{Constitutional-political vs party-political secularism:}
\begin{itemize}
\item \textbf{Constitutional-political secularism} --- equal respect for all religions; but every aspect cannot be respected, so the state may intervene when religious groups promote communal conflict or discriminate; personal laws may be diluted keeping the constitutional idea of social justice supreme. It is based on \textbf{principled distance}, a response to inter- and intra-religious domination.
\item \textbf{Party-political secularism} --- based not on principled distance but \textbf{opportunist distance} (affiliation/alliance with religious communities for electoral victory and mutual benefit).
\end{itemize}
\end{KeyConcept}

\begin{ExampleBox}
\textbf{Examples of opportunist distance:}
\begin{itemize}
\item The ban on \textit{The Satanic Verses}; unlocking the Babri Masjid for Ram Temple pooja rituals; withdrawal of the Haj subsidy; a state law preventing Hindu--Muslim marriage (framed as preventing ``forced conversion'') --- all opportunist distance.
\end{itemize}
\end{ExampleBox}

\begin{CriticalPoint}
\textbf{``Untrammeled majoritarianism masquerading as secularism'':}
\begin{itemize}
\item Bhargava's one-line summary: \textbf{untrammelled majoritarianism masquerading as secularism} = party-political secularism. The very idea of Indian secularism is to ensure the Hindu-majority state does not become a Hindu-majoritarian state.
\item A \textbf{feminist} reading of the same intervention (e.g., UP's ``love-jihad'' law) would say: by intervening in the private affair of marriage ``to protect women'', the state assumes women are not free --- that is a \textbf{patriarchal} state (Amartya Sen: it's about being a ``maternalist'' state).
\end{itemize}
\end{CriticalPoint}

\begin{QuickRevision}
\begin{itemize}
\item Constitutional-political = principled distance (equal respect + intervention for social justice). Party-political = opportunist distance (electoral alliances) = ``untrammelled majoritarianism masquerading as secularism.''
\item Feminist reading: state protecting women assumes women aren't free $\rightarrow$ patriarchal.
\end{itemize}
\end{QuickRevision}

\sectiondivider

\subsection{Amartya Sen: Symmetrical Treatment}

\begin{KeyConcept}
\begin{itemize}
\item Secularism is neither principled nor opportunist distance --- it is \textbf{symmetrical treatment}. If you treat everyone symmetrically, you are secular (even a theocratic state can be secular by this standard).
\end{itemize}
\end{KeyConcept}

\begin{ExampleBox}
\begin{itemize}
\item \textbf{Cow slaughter} is banned (the cow is a totem of one community), but pork slaughter is not banned (also a totem of another community) --- unequal, majoritarian. Symmetrical treatment = prohibit killing of \textit{all} totem animals.
\item In the \textbf{Indian army}, a Sikh has a Sikh regiment and may keep a turban, but a Muslim has no Muslim regiment and cannot keep a long beard --- not equal respect. Symmetrical treatment would treat all equally.
\item (Note: Sen's ``symmetry'' has not been fully detailed by him; an EPW article notes he did not clearly spell out its forms. A theocracy imposing one community's norms on all is also ``symmetrical'' at the surface, but may violate universal human rights.)
\end{itemize}
\end{ExampleBox}

\begin{DataFact}
\textbf{Conclusion:}
\begin{itemize}
\item Secularism in India has had a \textbf{tumultuous ride} --- imposed from above and developed from within; the conflict between the two is what we see in everyday life.
\end{itemize}
\end{DataFact}

\begin{QuickRevision}
\begin{itemize}
\item Sen: secularism = symmetrical treatment (cow vs pork slaughter; Sikh turban vs Muslim beard).
\item Conclusion: secularism = imposed from above vs developed from within (tumultuous ride).
\end{itemize}
\end{QuickRevision}

\sectiondivider

\subsection{Pluralism}

\begin{KeyConcept}
\textbf{Meaning:}
\begin{itemize}
\item \textbf{Pluralism} = peaceful, harmonious \textbf{co-existence of multiple world-views and belief systems} (not just religions). Multiple religions coexisting is only \textit{one} world-view; we also have coexistence of science and religion, of secularism and religion, and of the state, religion and science together. \textbf{Pluralism is a prerequisite for secularism} --- you cannot maintain state--religion distance unless there is plurality.
\item Its roots lie in \textbf{Akbar, Ashoka, Gandhi, Kabir} (who spoke of ``religious tolerance'', not ``secularism'').
\end{itemize}
\end{KeyConcept}

\begin{KeyConcept}
\textbf{Methodological pluralism:}
\begin{itemize}
\item Within sociology, Marxism, functionalism, feminism and postmodernism coexist under one umbrella --- \textbf{methodological pluralism}.
\end{itemize}
\end{KeyConcept}

\begin{DataFact}
\textbf{Ashoka's trajectory \& Freud's ``collective neurosis'' (student doubts):}
\begin{itemize}
\item \textbf{Ashoka} was initially intolerant, then became increasingly tolerant of all religions --- from ``Dheri Ghosh'' to ``Dham Ghosh''; cultural conquest, no more physical conquest. (In today's language we would call this secularism.)
\item \textbf{Freud and Marx} are the only two who have talked exclusively about the \textit{dysfunctions} of religion (besides the feminists, who follow Marx). \textbf{Freud}: religion is \textbf{wish-fulfilment} --- as primates we developed sexual attraction to the mother and collectively killed the father; this guilt (millions of years old) is transmitted across generations and surfaces as unconscious worship --- a form of \textbf{neurosis/OCD}.
\end{itemize}
\end{DataFact}

\begin{QuickRevision}
\begin{itemize}
\item Pluralism = coexistence of multiple world-views (not just religions); a prerequisite for secularism.
\item Roots: Akbar, Ashoka, Gandhi, Kabir (tolerance). Methodological pluralism = many paradigms in sociology.
\end{itemize}
\end{QuickRevision}

\sectiondivider

\section{Pluralism, Science \& Religion, and Religious Revivalism}

\subsection{Nehruvian Secularism \& Ashis Nandy}

\begin{CriticalPoint}
\textbf{The three Indian scholars are the Marx-Weber-Durkheim of Indian sociology:}
\begin{itemize}
\item Ashis Nandy, T. N. Madan and Partha Chatterjee are \textbf{more important than Marx, Weber and Durkheim} for Indian sociology --- they recur in every topic.
\end{itemize}
\end{CriticalPoint}

\begin{DataFact}
\textbf{Nehruvian secularism:}
\begin{itemize}
\item The secularism we practise is \textbf{Western secularism} (modified, but Western in origin) --- borrowed by \textbf{Nehru} (``Nehruvian secularism''). It works on \textbf{strict separation of religion from public life} (religion = private affair). Example: \textbf{Charlie Hebdo} published a cartoon of Prophet Muhammad; because religion was private, the state could not resolve it, and religion bounced back as \textbf{fundamentalism}.
\item Nehruvian secularism is part of the larger Western package of \textbf{nation-building, scientific temper and developmentalism} (nationalism from the French Revolution; development from the West/industrial revolution).
\item (Note: \textbf{Turkish secularism}, established by \textbf{Atat\"urk}, is also inspired by strict public--private separation, borrowed from the West/democracy.)
\end{itemize}
\end{DataFact}

\begin{CriticalPoint}
\textbf{Ashis Nandy on the result:}
\begin{itemize}
\item Applying Western secularism to India \textbf{weakens religious faith} (nation-building: ``we are Indians, not Hindu-Muslim''), but \textbf{strengthens religion as political-communalist ideology} (Citizenship Amendment Act, Article 370 --- justified as nation-building).
\item Its \textbf{latent function}: it establishes the Indian state as \textbf{legitimate arbiter among religious communities} --- whatever the state says is seen as non-discriminatory and justified ``in the name of secularism'', even when it discriminates (e.g., UP anti-conversion laws in the name of individual rights).
\item But this Western secularism \textbf{could not prevent communal conflicts}, nor promote tolerance --- hence we should return to the indigenous \textbf{philosophy of tolerance} (Hinduism, Islam, Buddhism --- Ashoka, Akbar, Gandhi).
\end{itemize}
\end{CriticalPoint}

\begin{QuickRevision}
\begin{itemize}
\item Nehruvian secularism = Western, separation of religion from public life (Charlie Hebdo example).
\item Nandy: it weakens faith but strengthens religion-as-ideology; latent function = state as legitimate arbiter; it failed to prevent communal conflict $\rightarrow$ return to indigenous tolerance.
\end{itemize}
\end{QuickRevision}

\sectiondivider

\subsection{Pluralism \& Peter Berger}

\begin{KeyConcept}
\textbf{Pluralism (recap):}
\begin{itemize}
\item Pluralism = mutual peaceful coexistence of multiple religious faiths (religious pluralism); also the coexistence of science and religion (scientific in public, religious in private). \textbf{Pluralism is a prerequisite for secularism.}
\end{itemize}
\end{KeyConcept}

\begin{CriticalPoint}
\textbf{Peter Berger:}
\begin{itemize}
\item \textbf{Peter Berger} (who proposed the secularization thesis and later refuted it): \textbf{modernity is characterised not by secularity but by plurality}. We may not find secularism in its true form, but we find \textbf{pluralism} (people coexist peacefully even if the state is not equally distant) --- a product of modern education. America practises secularism \textit{within} its territory and fundamentalism \textit{outside} (on a mission to convert theological states into democracies).
\end{itemize}
\end{CriticalPoint}

\begin{CriticalPoint}
\textbf{Threats to pluralism:}
\begin{enumerate}
\item \textbf{Development} --- fruits of development are appropriated by dominant castes (Reddy, Yadav, Patidar, Jat); when those left behind demand reservation, it creates conflict (Sachar Committee: Muslims lag behind Hindus; $\sim$20\% upper-caste men hold $\sim$60\% of higher-education opportunities --- Satish Deshpande's \textbf{exclusive inequality}).
\item \textbf{Boundary maintenance} --- ghettoization of minorities (Muslim ghettos).
\item \textbf{Fundamentalism, revivalism, communalism, communalization of politics, hate speech, laws that promote conflict} (e.g., anti-conversion laws).
\end{enumerate}
\end{CriticalPoint}

\begin{QuickRevision}
\begin{itemize}
\item Berger: modernity = plurality, not secularity (he refuted his own secularization thesis).
\item Threats to pluralism: development (exclusive inequality), boundary maintenance, fundamentalism/communalism/hate speech/anti-conversion laws.
\end{itemize}
\end{QuickRevision}

\sectiondivider

\subsection{Religion and Science: Three Models}

\begin{KeyConcept}
\textbf{What is at stake:}
\begin{itemize}
\item Religion and science are not temple vs rocket --- they are \textbf{bodies of knowledge}. The conflict between them gave birth to \textbf{sociology} (the child of the conflict between these two ideologies). Science = ``with tons of evidence, I still doubt what I see'' (critical); religion = ``without evidence, I believe what I don't see.'' Both are creations of the human mind.
\end{itemize}
\end{KeyConcept}

\begin{DataFact}
\textbf{The three models:}
\begin{enumerate}
\item \textbf{Evolutionary model} (Comte --- theological/metaphysical/scientific; Frazer --- magic$\rightarrow$religion$\rightarrow$science; Weber --- rationalization/disenchantment) --- science replaces religion.
\item \textbf{NOMA (Non-Overlapping Magisteria)} --- \textbf{Stephen Jay Gould}: science and religion occupy separate, non-overlapping domains (different houses in one city). Science answers ``how'' (laws, gravitation, Big Bang), religion answers ``why'' (why the laws exist, why the Big Bang) --- a division of labour/organic solidarity. (Every great scientist --- Hawking, Einstein, Newton, Dawkins --- was also a believer; ``Eism'', religion of the enlightened.)
\item \textbf{Fusion model} --- a critique of NOMA: religion and science \textbf{overlap} and help each other.
\end{enumerate}
\end{DataFact}

\begin{ExampleBox}
\textbf{Examples mapped to models:}
\begin{itemize}
\item Religion for uncertain / science for certain $\rightarrow$ \textbf{NOMA} (Malinowski). NASA eating ``lucky peanuts'' before a launch; ISRO director at Tirupati before Chandrayaan; CERN naming the ``\textbf{God particle}''; Scientology and UFO religions (religion built on scientific machines); farmers' rain rituals $\rightarrow$ \textbf{fusion}. An atheist scientist marrying with rituals / an Indian wedding (rituals as religion) $\rightarrow$ \textbf{NOMA}. Amniocentesis restricted by ethics; ART/three-parent baby banned (ethics from religion govern science); testing the Gayatri mantra's efficacy on COVID $\rightarrow$ \textbf{fusion}.
\end{itemize}
\end{ExampleBox}

\begin{CriticalPoint}
\textbf{Conclusion \& critics:}
\begin{itemize}
\item Conclusion: \textbf{there is no legitimate conflict between science and religion} (Einstein: ``science without religion is lame; religion without science is blind''). The relation is \textbf{dynamic} --- once seen as conflict, today as cooperation.
\item \textbf{Karl Popper}: religion is a \textbf{closed system} (not falsifiable), science is an \textbf{open system} (provisional, open to criticism). \textbf{Thomas Kuhn}: both are \textbf{equally dogmatic} --- science is conservative (it takes long to accept a paradigm shift, e.g., Newton$\rightarrow$Einstein); religion is also open (religious reforms, women monks). \textbf{Paul Feyerabend}: both are the same (``anything goes'').
\item \textbf{Yuval Noah Harari} (\textit{Homo Deus}): technology (stem cells, DNA, CRISPR-Cas9) may make man himself God/immortal --- the scientific temper that Weber said arose from religion now makes man superhuman.
\end{itemize}
\end{CriticalPoint}

\begin{QuickRevision}
\begin{itemize}
\item Three models: evolutionary (science replaces religion), NOMA (separate domains, Gould), fusion (overlap and help each other).
\item Einstein: no legitimate conflict. Popper (religion closed/science open), Kuhn (both dogmatic), Feyerabend (both same). Harari: man becoming god.
\end{itemize}
\end{QuickRevision}

\sectiondivider

\subsection{Religious Revivalism}

\begin{KeyConcept}
\textbf{Meaning (multiple forms):}
\begin{itemize}
\item \textbf{Religious revivalism} = reviving/returning religion to life. It has \textbf{multiple forms} (no single meaning). Durkheim: religion forms around the sacred; ``God'' is the collective effervescence of gathered people (as in a football stadium) --- God doesn't exist, it is the excited state.
\end{itemize}
\end{KeyConcept}

\begin{DataFact}
\textbf{Forms of revivalism:}
\begin{enumerate}
\item \textbf{Revivalist movements} --- \textbf{Dayanand Saraswati's Shuddhi movement} (Ghar Wapsi --- bringing Hindu converts back); the Deoband movement; the Arya Samaj. (Historians split movements into ``reformist'' vs ``revivalist'' --- de Bonald/de Maistre as revivalist; but a person like Dayanand is \textit{both} reformer and revivalist --- you cannot reform without reviving. These are models, not absolute reality.)
\item \textbf{Protestant Reformation} (Martin Luther reinterpreted religion --- revival + reform).
\item \textbf{New world religions} --- proliferation of cults, spiritual revolution (crystal healing, meditation) --- spirituality is just a new form of religion (Stark \& Bainbridge); the ``spiritual, not religious'' claim is \textbf{role distance} (Goffman).
\item \textbf{Religious conversion} --- Ambedkar and the Hathras community converting to (Neo-)Buddhism to liberate themselves from caste.
\item \textbf{Individual revival} --- praying on exam results day, Sachin Tendulkar looking up at the sky after a century.
\end{enumerate}
\end{DataFact}

\begin{CriticalPoint}
\textbf{Revivalism challenges the secularization thesis:}
\begin{itemize}
\item Modernization/secularization thesis established the decline of religion; but the world has witnessed \textbf{revivalism many times}. \textbf{Stark \& Bainbridge}: perpetual cycle of religious revival and renewal. \textbf{Norris \& Inglehart} (existential security): poor/insecure societies revive religion more. \textbf{Durkheim}: as long as there is society, there is the sacred. \textbf{David Lyon}: private religions, re-enchantment.
\item Revivalism can \textbf{justify} the social order (revivalism) or \textbf{challenge} it (reformism). Science itself revives religion (the ``God particle''; the idea of immortality coming from religion, aided by stem cells).
\item \textbf{Marx}: the bourgeoisie revive religion to legitimize exploitation; \textbf{feminists}: patriarchy revives religion to legitimize gender inequality.
\end{itemize}
\end{CriticalPoint}

\begin{QuickRevision}
\begin{itemize}
\item Revivalism = reviving religion; multiple forms (Shuddhi/Ghar Wapsi, Protestant Reformation, new cults/spirituality, conversion, individual revival).
\item It challenges secularization (perpetual cycle --- Stark \& Bainbridge; existential security --- Norris \& Inglehart; Durkheim's sacred; re-enchantment --- David Lyon). Marx/feminists: elites revive religion to justify domination.
\end{itemize}
\end{QuickRevision}

\sectiondivider

\section{Fundamentalism}

\subsection{Answer-Writing: Corruption via Manifest \& Latent Analysis}

\begin{PYQLink}
\begin{itemize}
\item ``Using manifest and latent analysis, discuss the prevalence of corruption in Indian society.'' (Asked in 2014, reflecting the 2012 anti-corruption movement.)
\end{itemize}
\end{PYQLink}

\begin{ExampleBox}
\textbf{Manifest (anticipated) functions of corruption:}
\begin{itemize}
\item Loss of public trust in the state; reduced efficiency/rationalization; promotes further corruption; drainage of public resources; breach of rules/code of conduct; abuse of power; favouritism.
\end{itemize}
\end{ExampleBox}

\begin{ExampleBox}
\textbf{Latent (unanticipated) functions of corruption:}
\begin{itemize}
\item Rise of the \textbf{anti-corruption movement} (2012, all-India); change in power structure (Congress $\rightarrow$ BJP); rise of \textbf{charismatic leaders} (Arvind Kejriwal --- the movement died but charisma routinized into the Aam Aadmi Party); rise of \textbf{civil society} (Transparency International gained visibility --- corruption became a measure of state efficiency); surveillance (Lokpal); tax havens; demoralization of honest public servants; an apathetic attitude toward bureaucracy (the policeman seen as a bribe-taker); reinforces materialism (Marx's M1--M2: money becomes the end).
\end{itemize}
\end{ExampleBox}

\begin{QuickRevision}
\begin{itemize}
\item Manifest: loss of trust, inefficiency, resource drain, abuse of power. Latent: anti-corruption movement (2012), Congress$\rightarrow$BJP, Kejriwal/charisma$\rightarrow$AAP, Transparency International, Lokpal/surveillance, demoralization.
\end{itemize}
\end{QuickRevision}

\sectiondivider

\subsection{Fundamentalism: Definition \& Features}

\begin{KeyConcept}
\textbf{Definition:}
\begin{itemize}
\item \textbf{Fundamentalism} = \textbf{unwavering commitment to an extreme ideology} (not just religion). If you are not open to another version of reality, you are a fundamentalist. If the ideology is religious belief, it is \textbf{religious fundamentalism}. (The word ``fundamental'' comes from Christianity --- to stick to the fundamentals of the religious text.)
\item Examples beyond religion: \textbf{gender/sexual fundamentalism} (rejecting a third gender), \textbf{political fundamentalism} (``if you oppose the leader you are anti-national''), \textbf{scientific fundamentalism} (Newton over Einstein --- Kuhn's conservatism). Media has made us equate fundamentalism with ISIS/Al-Qaeda (cultural hegemony).
\end{itemize}
\end{KeyConcept}

\begin{DataFact}
\textbf{The seven features (memorise):}
\begin{enumerate}
\item \textbf{Literal interpretation of the text} (Brahmin born of Brahma's mouth; Jesus born to the Virgin Mary; Shariat law).
\item \textbf{Charismatic leadership} (Osama bin Laden, Dayanand Saraswati, Veer Savarkar, Golwalkar).
\item \textbf{Rejection of religious pluralism} (only my religion has absolute truth; others are ignorant --- Dayanand: Hinduism is the only true religion).
\item \textbf{Self-appointed guardian of religion} (modernity has corrupted people --- pornography, abortion, transgender rights --- so they must restore religion's dignity).
\item \textbf{Rejection of modern reproductive technology} (surrogacy, three-parent baby --- commercialization of the womb).
\item \textbf{Use of modern technology} (TV, internet, WhatsApp, modern weapons --- against modernity but using its tools).
\item \textbf{Opposition to modernity} --- highly conservative; proposes to change contemporary liberal-pluralist society back into a conservative, traditional one. (Also: \textbf{legitimizes patriarchy}.)
\end{enumerate}
\end{DataFact}

\begin{QuickRevision}
\begin{itemize}
\item Fundamentalism = unwavering commitment to an extreme ideology; religious fundamentalism = one form.
\item Seven features: literal text, charismatic leadership, reject pluralism, guardian of religion, reject ART, use modern technology, oppose modernity (+ legitimize patriarchy).
\end{itemize}
\end{QuickRevision}

\sectiondivider

\subsection{Causes of Fundamentalism}

\begin{CriticalPoint}
\textbf{The causes:}
\begin{enumerate}
\item \textbf{Response to modernization, secularization and cosmopolitanism} (Steve Bruce, Anthony Giddens) --- push religion out the front door and it returns through the back door. Modern ideas (abortion, pornography, birth control, women's/homosexual rights, divorce, live-in relationships, euthanasia, female genital mutilation) undermine traditional beliefs and are seen as morally corrupt. (WHO: $\sim$200 million girls victims of FGM.)
\item \textbf{Monolithic/Abrahamic religions} --- religions with one text (Islam's Quran, Christianity's Bible, Sikhism's one holy book --- Khalsa) more likely to develop literal interpretation; religions with multiple texts (Hinduism, Buddhism's Tripitaka) allow multiple interpretations.
\item \textbf{Common enemy} --- ideological cohesion of hatred against one group (the Middle East vs America/Israel) channelizes fundamentalism.
\item \textbf{Western imperialism} --- US/UK intervention in the Middle East has the unintended consequence of fundamentalist groups.
\item \textbf{Globalization} --- spreads ideas (and fundamentalism) faster; Samuel Huntington's \textbf{clash of civilizations} (Western vs Eastern).
\end{enumerate}
\end{CriticalPoint}

\begin{ExampleBox}
\textbf{The Iranian Revolution:}
\begin{itemize}
\item The \textbf{Shah} (installed and backed by America in the 1930s, during the global anomie before WWII) introduced Western/modern ideas (clubs, bars, casinos, women's education/equality) to a traditional theocratic Iran. In 1979, \textbf{Ayatollah Khomeini} --- like de Bonald/de Maistre in France --- countered these changes, destroyed the modern lifestyle, and restored religion (back to the Quran). He was a \textbf{revivalist and a fundamentalist} --- and what differentiates a fundamentalist from a mere revivalist is the \textbf{fight for power} (ISIS wants to establish a caliphate).
\end{itemize}
\end{ExampleBox}

\begin{CriticalPoint}
\textbf{Fundamentalism as unanticipated consequence:}
\begin{itemize}
\item Fundamentalism is an \textbf{unanticipated consequence of modernization and secularization}. Modernity came as a response to religion; religion revives as a response to modernity (a cycle).
\end{itemize}
\end{CriticalPoint}

\begin{QuickRevision}
\begin{itemize}
\item Causes: modernization/secularization/cosmopolitanism (Bruce, Giddens), monolithic religions, common enemy, Western imperialism, globalization (Huntington).
\item Iranian Revolution (Shah vs Khomeini, 1979); fundamentalism = revivalism + fight for power; it is an unanticipated consequence of modernity.
\end{itemize}
\end{QuickRevision}

\sectiondivider

\subsection{Secular Fundamentalism}

\begin{CriticalPoint}
\textbf{``Secular fundamentalism'' (Grace Davie):}
\begin{itemize}
\item \textbf{Grace Davie} coined \textbf{secular fundamentalism}: being rigidly committed to an idea that \textit{seems} modern (e.g., ultra-nationalism --- ``if you don't obey, you are anti-national''). It is the same tendency as religious fundamentalism --- there is no difference between \textbf{heresy} (old church) and \textbf{contempt of court} (modern court): the modern court is the secular avatar of the church.
\end{itemize}
\end{CriticalPoint}

\begin{CriticalPoint}
\textbf{Surveillance \& docile bodies (Michel Foucault):}
\begin{itemize}
\item The \textbf{latent function of surveillance} (Pegasus, CCTV) is to produce \textbf{docile bodies} --- bodies that conform without being told to change (you stop at a red light even without a policeman). Modern democracies have become surveillance societies $\rightarrow$ gradually no difference between democracy, autocracy and theocracy --- the \textbf{end of ideology} (Daniel Bell).
\end{itemize}
\end{CriticalPoint}

\begin{QuickRevision}
\begin{itemize}
\item Secular fundamentalism (Davie) = rigid commitment to ``modern'' ideas (ultra-nationalism); contempt of court = secular heresy.
\item Foucault: surveillance $\rightarrow$ docile bodies $\rightarrow$ surveillance societies $\rightarrow$ end of ideology (Daniel Bell).
\end{itemize}
\end{QuickRevision}

\sectiondivider

\section{Politics \& Society --- Theories of Power}

\subsection{Answer-Writing: Religious Beliefs in Pre-Modern Societies}

\begin{PYQLink}
\begin{itemize}
\item ``State the reasons for the various religious beliefs in pre-modern societies.'' (2020, 10 marks.) Keyword = religious beliefs; context = pre-modern.
\end{itemize}
\end{PYQLink}

\begin{DataFact}
\textbf{The beliefs and their reasons:}
\begin{itemize}
\item \textbf{Totemism} (Durkheim, Arunta tribe) --- belief in the sacred/collective.
\item \textbf{Trobrianders} (Malinowski) --- religiosity from \textbf{uncertainty} (certain = science, uncertain = religion).
\item \textbf{Animism} (E. B. Tylor) --- belief in spirits (animate and inanimate), from the intellectual need to explain \textbf{death and dream}.
\item \textbf{Naturism} (Max M\"uller) --- worship of natural forces (fire/Agni, earth goddess, river), because nature in its extreme/calamitous form (flood, forest fire) brings crisis, so people worship it to avoid its wrath.
\item (Two further religions studied only by anthropologists: \textbf{animatism}.)
\end{itemize}
\end{DataFact}

\begin{QuickRevision}
\begin{itemize}
\item Totemism (sacred), Trobrianders (uncertainty), Animism (spirits; death/dream), Naturism (natural forces; crisis).
\item These beliefs persist even in modern societies.
\end{itemize}
\end{QuickRevision}

\sectiondivider

\subsection{Politics \& Society: Sociologist vs Political Scientist}

\begin{CriticalPoint}
\textbf{How a sociologist differs from a political scientist:}
\begin{itemize}
\item A political scientist classifies democracies by form (presidential, parliamentary, semi-presidential). A sociologist \textbf{questions the classification itself} --- e.g., \textbf{Andr\'e B\'eteille} and \textbf{Ambedkar}: democracy should be evaluated by the \textbf{strength/functioning of the opposition}, not the form of government. (Organizations like the Economist Intelligence Unit rate democracies as ``flawed''.)
\item Political scientists look only at the surface (they don't study religion's role in politics); sociologists study everything and reveal the hidden.
\item \textbf{Power is to society what blood is to the body} --- it flows everywhere; you need a lens to see it. Politics happens even in families (property disputes).
\end{itemize}
\end{CriticalPoint}

\begin{DataFact}
\textbf{The seven theories of power:}
\begin{enumerate}
\item Functionalist
\item Marxist
\item Weberian
\item Neo-Marxist
\item Pluralist
\item Elite
\item Feminist
\item Postmodern
\end{enumerate}
\end{DataFact}

\begin{QuickRevision}
\begin{itemize}
\item Sociologist questions the classification (B\'eteille/Ambedkar: opposition as the test of democracy); power = blood of society.
\item Eight theories: functionalist, Marxist, Weberian, neo-Marxist, pluralist, elite, feminist, postmodern.
\end{itemize}
\end{QuickRevision}

\sectiondivider

\subsection{Functionalist Theory of Power}

\begin{KeyConcept}
\textbf{Parsons: power as a resource:}
\begin{itemize}
\item \textbf{Talcott Parsons}: power is a \textbf{resource} used to help society meet its needs. Like depositing money in a bank (getting interest), we deposit our power in a leader, who returns it as \textbf{social and economic welfare} (roads, electricity, education, subsidies, oxygen/ventilators). The \textbf{system benefits} --- everybody wins. If the leader fails, the system deteriorates and everybody loses.
\item Therefore power is \textbf{variable-sum} (everybody wins or everybody loses).
\end{itemize}
\end{KeyConcept}

\begin{QuickRevision}
\begin{itemize}
\item Parsons: power = resource (deposit $\rightarrow$ interest = welfare); system benefits; variable-sum.
\end{itemize}
\end{QuickRevision}

\sectiondivider

\subsection{Marxist Theory of Power}

\begin{KeyConcept}
\textbf{Power in the economic base:}
\begin{itemize}
\item To Marx, power \textbf{caters to the bourgeoisie} who exploit the working class. Power comes from the \textbf{economic base} --- whoever controls the economy controls society; politicians are puppets of corporates (state is superstructure).
\end{itemize}
\end{KeyConcept}

\begin{ExampleBox}
\begin{itemize}
\item Commonwealth Games: functionalists see everyone benefiting; Marx sees Japan (and the Japan--India corporate alliance) making the profit --- the celebration of ``Metroman'' E. Sreedharan hides who profited. Demonetization benefited ATM/Paytm banks (new consumers), not the masses without smartphones (``Paytm'' became common sense). Rural electrification, urban housing, e-RUPI --- all cater to corporates.
\end{itemize}
\end{ExampleBox}

\begin{CriticalPoint}
\textbf{Power is constant-sum:}
\begin{itemize}
\item Power is \textbf{constant-sum} (+2/$-$2 = 0): one group's gain is another's loss (profit = exploitation). There is no collective gain or loss --- only corporate gain. Hence ``workers of the world unite'' $\rightarrow$ \textbf{dictatorship of the proletariat}.
\end{itemize}
\end{CriticalPoint}

\begin{QuickRevision}
\begin{itemize}
\item Marx: power in economic base; corporates/politicians puppets; constant-sum (profit = exploitation); dictatorship of the proletariat.
\end{itemize}
\end{QuickRevision}

\sectiondivider

\subsection{Weberian Theory of Power}

\begin{KeyConcept}
\textbf{Power as chance:}
\begin{itemize}
\item \textbf{Weber}: power is the \textbf{probability/chance} of a man or a number of men to exercise their will even against the resistance of others. The word \textbf{chance} is crucial --- it is Weber's reply to Marx's determinism: power can be held by \textit{anyone} (father in the family, teacher, monitor), not just owners of the means of production.
\item Exercising will against resistance = domination = one's gain is another's loss (so power is again \textbf{constant-sum}). Authority: traditional, charismatic, legal-rational.
\end{itemize}
\end{KeyConcept}

\begin{QuickRevision}
\begin{itemize}
\item Weber: power = chance to exercise will against resistance (anti-determinism); constant-sum; three types of authority.
\end{itemize}
\end{QuickRevision}

\sectiondivider

\subsection{Neo-Marxist Theories: Althusser \& Gramsci}

\begin{KeyConcept}
\textbf{Althusser: superstructure is relatively autonomous:}
\begin{itemize}
\item \textbf{Louis Althusser}: the superstructure (forms of consciousness) is \textbf{relatively autonomous} from the economic base --- you can be powerful by controlling the \textit{superstructure}, not just the means of production.
\item \textbf{Two apparatuses (tools) of the state:}
\begin{enumerate}
\item \textbf{Ideological State Apparatus (ISA)} --- state controls ideology (religion, family, law, education, media). State plans families (ration cards, IDs, contraceptives); state and religion are ``friends with benefits''; state controls education (the first \textbf{Saraswati Shishu Mandir} was established in 1952; \textbf{saffronization of education} --- reading Savarkar/Golwalkar instead of Gandhi; teaching Muslims as invaders --- \textbf{communalization of education}; a 2017 NCERT case against the Shishu Mandir curriculum for being against India's plural-democratic fabric).
\begin{itemize}
\item The NCERT itself makes \textbf{Gandhi a super-hero} while \textbf{Bhagat Singh, Subhas Chandra Bose and Sardar Patel} are marginalised (the radical is always projected as unwanted) --- because the post-independence curriculum was designed by the Congress party.
\end{itemize}
\item \textbf{Repressive State Apparatus (RSA)} --- police, army, judiciary, which act for the state against those who disobey (JNU/Delhi University protests).
\end{enumerate}
\item Althusser gives \textbf{no agency to individuals} --- \textbf{structural Marxism} (individuals are conformist, manipulated).
\item (Context: Althusser is French, wrote against students' protests of the 1950s/60s; Gramsci is Italian, wrote against Mussolini's fascism.)
\end{itemize}
\end{KeyConcept}

\begin{KeyConcept}
\textbf{Gramsci: humanist Marxism:}
\begin{itemize}
\item \textbf{Gramsci} also sees the political society commanding the civil society (religion, education, media) via \textbf{cultural hegemony} --- but he \textbf{gives agency to individuals}: organic intellectuals can produce \textbf{counter-hegemony}. Hence his version is \textbf{humanist Marxism} (vs Althusser's structural Marxism).
\end{itemize}
\end{KeyConcept}

\begin{QuickRevision}
\begin{itemize}
\item Althusser: superstructure relatively autonomous; ISA (religion/family/law/education/media) + RSA (police/army/judiciary); structural Marxism (no agency).
\item Gramsci: civil + political society, cultural hegemony + counter-hegemony $\rightarrow$ humanist Marxism.
\end{itemize}
\end{QuickRevision}

\sectiondivider

\subsection{Pluralist Theories: Dahl \& Marsh}

\begin{KeyConcept}
\textbf{Robert Dahl (classical pluralism):}
\begin{itemize}
\item \textbf{Robert Dahl} (\textit{Who Governs?}) critiqued Parsons: there is no ``system'' that is powerful --- power lies with \textbf{interest groups} (pressure groups; even opposition parties, student unions, teacher associations, IAS Officers Association). Society consists of \textbf{multiple interest groups, each equally influential}, so \textbf{power is dispersed} among them.
\item When a state formulates policy (lateral entry; farm loan waivers vs bank NPAs; 100\% FDI in defence), multiple stakeholders make conflicting demands; the state takes the interest of \textbf{everyone} (hence bills are publicly uploaded for feedback before becoming acts).
\end{itemize}
\end{KeyConcept}

\begin{KeyConcept}
\textbf{David Marsh (elite pluralism):}
\begin{itemize}
\item \textbf{David Marsh} critiqued Dahl: all interest groups are \textbf{not} equally influential --- \textbf{some (elite) interest groups are more powerful} (wealth, influential leaders, educated urban upper-caste men). Not everyone's interest is taken --- the interest of a \textbf{few elites} is taken, the rest rejected. The government is \textbf{not an honest broker}; \textbf{``governance is the business of compromise''} (Raymond Aron).
\end{itemize}
\end{KeyConcept}

\begin{QuickRevision}
\begin{itemize}
\item Dahl: power dispersed among equally influential interest groups (critique of Parsons).
\item Marsh: some (elite) interest groups dominate; ``governance is the business of compromise.''
\end{itemize}
\end{QuickRevision}

\sectiondivider

\section{Elite Theory \& Postmodern Theory}

\subsection{Answer-Writing: ``World Religions Patriarchal''}

\begin{PYQLink}
\begin{itemize}
\item ``World religions are patriarchal. Substantiate your answer with examples.'' (20 marks, but treat as 10 for point-building.)
\end{itemize}
\end{PYQLink}

\begin{KeyConcept}
\textbf{The structure:}
\begin{itemize}
\item No theory is absolutely right. The question really asks: \textbf{is the feminist theory of religion right?} First argue \textit{why} religions are seen as patriarchal (Simone de Beauvoir --- women as second-class citizens, consciousness nullified, promised a better hereafter, serving patriarchy), citing the 4--5 parameters (religious texts, religious organisations, leadership, role in everyday life, food/dress code, laws). Then present the \textbf{counter}: ``world religion'' is broader than mainstream religions --- new world religions (cults, spiritual revolution), animism/totemism/naturism (found everywhere), a Shakta sect purely of goddesses, women-run cults, and women monks (Sekar Babu, Tamil Nadu --- a reform prompted by the \textbf{Sabarimala controversy}). Sikhism is the \textit{least} patriarchal (but still patriarchal). Present both sides; conclude that ``patriarchal'' is subjective and highly open to interpretation.
\end{itemize}
\end{KeyConcept}

\begin{QuickRevision}
\begin{itemize}
\item Argue feminist case (Beauvoir; texts/organisations/leadership/laws) $\rightarrow$ counter (new world religions, cults, Shakta, women monks) $\rightarrow$ ``patriarchal'' is subjective.
\end{itemize}
\end{QuickRevision}

\sectiondivider

\subsection{Elite Theory: Pareto}

\begin{KeyConcept}
\textbf{Elite:}
\begin{itemize}
\item \textbf{Elite} = anyone who is \textbf{exceptional and ranks high in all the indices of their field} (Dhoni, Tendulkar, Ronaldo, M. F. Husain, Leonardo, Amitabh Bachchan, Amartya Sen, Abhijit Banerjee, Zakir Hussain; elite lawyers like Ram Jethmalani, Kapil Sibal, Harish Salve, Prashant Bhushan; elites even among criminals --- Dawood Ibrahim). Elites are found in \textbf{every} field and social identity (Dalit elites --- Ambedkar, Jyotiba Phule, Mayawati).
\end{itemize}
\end{KeyConcept}

\begin{DataFact}
\textbf{Pareto's typology \& circulation:}
\begin{itemize}
\item \textbf{Vilfredo Pareto} (founder of Italian sociology, critic of Marx, influenced by \textbf{Machiavelli's} \textit{The Prince}): the \textbf{birth of elites is inevitable}; the majority never rules --- the minority rules. Two types of elite:
\begin{itemize}
\item \textbf{Governing elite} = lions (aristocracy, speculators) --- rule by \textbf{force and coercion}, decisive and brave (military dictatorship, Hitler).
\item \textbf{Non-governing elite} = foxes (rentiers, working class) --- rule by \textbf{deceit and manipulation}, cunning.
\end{itemize}
\item \textbf{Circulation of elites}: governing elites decay over time and are replaced by non-governing elites (lions replaced by foxes). Lions maintain a \textbf{closed system} (no outsider; dynasty politics). The British came as foxes (economic adventurers) and became the new lions; Bahadur Shah Zafar was a decayed lion. \textbf{History is the graveyard of aristocracy} --- the history of all societies is the struggle between circulating elites, not class struggle.
\item The \textbf{ideal ruler} = a combination of \textbf{lion (bravery) and fox (cunning)} --- Machiavelli's prince. (Pareto and \textbf{Mosca} give a psychological theory of elite --- elites are born with special psychological traits.)
\end{itemize}
\end{DataFact}

\begin{CriticalPoint}
\textbf{Against Marx:}
\begin{itemize}
\item Even in communism, elites will arise (Mao Zedong; Nazism/fascism promised heaven but elites hoarded power). Democracy is a chamber of aristocracy.
\item Example: \textbf{Congress} (formed 1885, one of the world's oldest parties, older than the British Labour Party) has always been ruled by elites --- its presidents come from the Gandhi family or ruling dynasties (Scindia, Gaekwad, Chauhan) --- \textbf{dynasty politics} = rule of elites. ``Rule of the majority, by the majority, for the majority'' is never the reality; it is rule of the minority, by the minority, for the minority.
\end{itemize}
\end{CriticalPoint}

\begin{QuickRevision}
\begin{itemize}
\item Elite = exceptional, ranks high in a field. Pareto: birth of elite inevitable; lions (force) vs foxes (cunning); circulation of elite; history = graveyard of aristocracy; ideal ruler = lion + fox.
\item Against Marx: elites rule even in communism; ``rule of the minority for the minority.''
\end{itemize}
\end{QuickRevision}

\sectiondivider

\subsection{Power Elite: C. Wright Mills}

\begin{KeyConcept}
\textbf{The power elite:}
\begin{itemize}
\item \textbf{C. Wright Mills} (who gave sociological imagination and white-collar alienation): the \textbf{power elite} = the \textbf{heads of the most dominant organisations/institutions}. In America, three: the \textbf{military} (Navy SEALs = Sea, Air, Land), the \textbf{executive (federal) branch}, and the \textbf{corporate rich}. It is an \textbf{institutional} theory --- elites are those who command these positions (no non-governing elite).
\end{itemize}
\end{KeyConcept}

\begin{DataFact}
\textbf{Traits of the power elite:}
\begin{enumerate}
\item \textbf{Similar/homogenous background} --- white, upper class, Protestant, $\sim$90\% from Ivy League universities (Yale, Harvard, Cambridge, Stanford --- every president from Obama to Trump to the Bush/Koch families; Obama the only non-white).
\item \textbf{Closed group / elite self-recruitment} --- they fill their own vacancies (sons/daughters marry within the circle).
\item \textbf{Their interests coincide} --- they work to secure each other's needs (a war benefits military promotion, corporate weapon sales, and presidential hegemony).
\end{enumerate}
\begin{itemize}
\item Mills calls the power elite a \textbf{quasi-hereditary caste} (endogamous, self-recruiting, born with silver spoons). External appearance is democracy; upper positions are closed.
\end{itemize}
\end{DataFact}

\begin{CriticalPoint}
\textbf{Context \& application:}
\begin{itemize}
\item Pareto wrote in the context of Nazism/fascism; Mills investigates American democracy. America governs the world (IMF, World Bank, WTO; the Kudankulam case --- Greenpeace vs Westinghouse, India favoured Westinghouse).
\end{itemize}
\end{CriticalPoint}

\begin{QuickRevision}
\begin{itemize}
\item Power elite (Mills) = heads of military + executive + corporate rich; homogenous (white, upper-class, Ivy League); elite self-recruitment; quasi-hereditary caste; institutional theory.
\end{itemize}
\end{QuickRevision}

\sectiondivider

\subsection{Postmodern Theory: Baudrillard}

\begin{KeyConcept}
\textbf{Simulacra \& hyperreality:}
\begin{itemize}
\item \textbf{Jean Baudrillard} (\textit{End of Politics}; postmodern books are often titled ``End of\dots'' or ``Death of\dots''): we have entered a world of \textbf{simulacra} --- signs that make us believe reality exists when it does not. Media produces \textbf{hyperreality}.
\item The \textbf{Gulf War}: before it he wrote ``the Gulf War will not happen''; during, ``is not happening''; after, ``never happened'' --- because we saw it only through media. \textbf{Balakot strike}: Pakistan says it never happened, India says it did. Today's ``war'' is a \textbf{game of war} (buildings crushed, few casualties, video-game visuals).
\item \textbf{Simulacra examples}: fake Nike shoes in Palika Bazaar (seem original); Fanta/Mirinda (more orange than orange, but not orange); TV debates (leaders appear opposed but are not); the films \textit{The Matrix} and \textit{Interstellar} (borrowed from the idea of simulacra).
\item Today we produce \textbf{signs}, not material goods. Tomorrow's world will be \textbf{media-saturated, hyper-real politics}.
\end{itemize}
\end{KeyConcept}

\begin{QuickRevision}
\begin{itemize}
\item Baudrillard: simulacra (signs that feign reality); hyperreality via media; Gulf War ``never happened''; fake Nike/Fanta/TV debates; Matrix/Interstellar; media-saturated hyper-real politics.
\end{itemize}
\end{QuickRevision}

\sectiondivider

\subsection{Postmodern Theory: Foucault}

\begin{KeyConcept}
\textbf{Sovereign power vs disciplinary power:}
\begin{itemize}
\item \textbf{Michel Foucault} is a critique of all theories of power (radical theorization). Two historical forms of power:
\begin{itemize}
\item \textbf{Sovereign power} (pre-modern): kings punished the body publicly (execution before an audience, boiling water) --- pain inflicted on the \textbf{body}.
\item \textbf{Disciplinary power} (modern, post-Enlightenment): punishment inflicted on the \textbf{soul}, not the body --- you discipline \textit{yourself} (a CCTV camera regulates your speed; you stop at a red light without a policeman; a camera in your room makes you self-conscious). This produces \textbf{docile bodies} --- bodies that conform without being told.
\end{itemize}
\item This is connected to \textbf{Frankfurt School Marxism} (media as superstructure) and \textbf{Gramsci} (civil society). Postmodernism is an extension of --- and critique of --- these theories.
\end{itemize}
\end{KeyConcept}

\begin{QuickRevision}
\begin{itemize}
\item Foucault: sovereign power (punish the body, pre-modern) $\rightarrow$ disciplinary power (punish the soul, self-regulation, docile bodies, surveillance/Pegasus); critique of all power theories.
\end{itemize}
\end{QuickRevision}